
base fidèlement plat quasi-compact. Si $i'$ est régulière, alors l'Idéal $\underline{J}$ qui définit $i$ est de type fini, et $\underline{J}/\underline{J}^2$ est localement libre de type fini, car ces propriétés se descendent par morphismes fidèlement plats quasi-compacts. Si $f_1,\ldots,f_d$ est une suite de sections de $\underline{J}$ relevant une base de $\underline{J}/\underline{J}^2$, le complexe de Koszul associé aux images inverses des $f_i$ est acyclique en degrés $>0$ d'après 1.4; comme c'est l'image inverse du complexe de Koszul associé aux $f_i$, ce dernier est acyclique en degrés $>0$, et $i$ est régulière.

\textbf{Définition 1.6.} \uline{Soit} $i:Y\longrightarrow X$ \uline{une immersion définie dans un ouvert} $U$ \uline{de} $X$ \uline{par un Idéal} $\underline{J}$ \uline{de} $\underline{O}_X$; \uline{le} $\underline{O}_Y$-\uline{Module} $\underline{J}/\underline{J}^2$ (\uline{indépendant de} $U$) \uline{est appelé} faisceau conormal de $Y$ dans $X$, \uline{et noté} $\underline{N}_{X/Y}$. \uline{Si} $i$ \uline{est régulière}, $\underline{J}/\underline{J}^2$ \uline{est localement libre de type fini}, \uline{et son rang s'appelle} codimension de $Y$ dans $X$, \uline{ou encore} \uline{codimension de} $i$.

Si on a deux immersions fermées : $Z\xrightarrow{j}Y\xrightarrow{i}X$, les faisceaux conormaux forment une suite exacte:
\begin{equation}\tag{1.6.1}
 j^*(\underline{N}_{X/Y})\longrightarrow \underline{N}_{X/Z}\longrightarrow \underline{N}_{Y/Z}\longrightarrow0.
\end{equation}
Soient en effet $\underline{I}$ l'idéal de $Y$ dans $X$, $\underline{J}$ l'idéal de $Z$ dans $X$, $\underline{K}$ l'idéal de $Z$ dans $Y$. On a la suite exacte:
\[
0\longrightarrow \underline{I}\longrightarrow \underline{J}\longrightarrow \underline{K}\longrightarrow0;
\]
d'où par tensorisation par $\underline{O}_X/\underline{I}=\underline{O}_Y$, la suite exacte:
\[
 \underline{I}/\underline{I}^2\longrightarrow \underline{J}\otimes_{\underline{O}_X}\underline{O}_Y\longrightarrow \underline{K}\otimes_{\underline{O}_X}\underline{O}_Y\longrightarrow0;
\]
en tensorisant par $\underline{O}_Z=\underline{O}_X/\underline{J}=\underline{O}_Y/\underline{K}$, on obtient la suite exacte:
\[
 j^*(\underline{I}/\underline{I}^2)\longrightarrow \underline{J}\otimes_{\underline{O}_X}\underline{O}_Z\longrightarrow \underline{K}\otimes_{\underline{O}_Y}\underline{O}_Z\longrightarrow0,
\]
qui est la suite cherchée.

